\section{Election, Baptism, and Sudden Consecration}

In the autumn of 374, Auxentius of Cappadocia died after nearly twenty years as Homoian bishop of Milan. Despite his condemnation by a Roman synod under Pope Damasus around 369, Auxentius had kept his see under the non-interventionist religious policy of Emperor Valentinian~I. His death reopened the unresolved wound of 355: Nicene Catholics and Homoians alike crowded the Milanese basilica to contest the succession, threatening open urban riot. As \latin{consularis} of Aemilia-Liguria, Ambrose entered the church to keep civic order and exhort the assembly to peace.

What happened next became one of the most celebrated reversals in late antiquity. The divided crowd turned from its rival candidates and acclaimed the unbaptized governor himself as bishop of Milan. Writing twelve years later to Emperor Valentinian~II during a renewed crisis over his see, Ambrose himself appealed to that popular election and to the elder Valentinian's ratification of it: ``I omit to mention that the people have themselves already given their judgment. I am silent as to the fact that they demanded of your father him whom they now have. I am silent as to the promise of your father that if he who was chosen would undertake the bishopric there should be tranquillity. I acted on the faith of these promises'' (\work{Letter} XXI.7).

Within months of the event, St.~Basil of Caesarea wrote from Cappadocia to congratulate the new bishop, confirming from the East how sudden and astonishing the transfer appeared to contemporaries:

\begin{quote}
I have given glory to God, Who in every generation selects those who are well-pleasing to Him; \ldots\ Who now has drawn forth for the care of Christ's flock a man from the imperial city, entrusted with the government of a whole nation, exalted in character, in lineage, in position, in eloquence, in all that this world admires. This same man has flung away all the advantages of the world, counting them all loss that he may gain Christ, and has taken in his hand the helm of the ship, great and famous for its faith in God, the Church of Christ. Come, then, O man of God; not from men have you received or been taught the Gospel of Christ; it is the Lord Himself who has transferred you from the judges of the earth to the throne of the Apostles; fight the good fight; heal the infirmity of the people, if any are infected by the disease of Arian madness; renew the ancient footprints of the Fathers. (St.~Basil the Great, \work{Letter} CXCVII.1)
\end{quote}

Ambrose's own admission in \work{De Officiis} I.1.2---that he could ``no longer now escape from the duty of teaching which the needs of the priesthood have laid upon us, though we tried to avoid it''---confirms his genuine reluctance. Later narrative in Paulinus of Milan's \work{Vita} (reported through H.~de Romestin's notes) fills in dramatic details: a child's sudden shout of ``Ambrose, bishop!'' inside the assembly, and Ambrose's desperate stratagems to disqualify himself by ordering judicial torture contrary to his custom and attempting flight by night before Emperor Valentinian~I confirmed the election. Those vivid particulars belong to fifth-century hagiographic memory and are examined separately below; the historical core established by Ambrose and Basil is already extraordinary enough.

Once Valentinian~I ratified the choice, Ambrose accepted it as the will of God, insisting on baptism from a Catholic rather than a Homoian hand. Eight days after his baptism, having passed rapidly through the minor and major clerical grades contrary to ordinary canonical intervals, he was consecrated bishop of Milan. Liturgical memory at Milan and Rome fixed the anniversary of his ordination on 7~December, which places his consecration on 7~December 374---a Sunday in 374, whereas modern scholars who note that 7~December fell on a Saturday in 373 debate whether the consecration occurred in 373 or 374. The weight of ancient and editorial consensus favors 374, while acknowledging that 373 remains a serious chronological alternative.

Ambrose began his episcopate with three decisive acts that set the pattern of the next twenty-three years. First, he stripped himself of personal control over his wealth. According to the early Milanese report summarized by de Romestin, he gave his gold and silver to the Church and the poor, reserved a life interest in his landed property for his sister Marcellina, and entrusted the entire stewardship of family estates to his brother Satyrus, who resigned his own public office to shield Ambrose from secular distractions. Ambrose's own words at Satyrus's funeral confirm this arrangement: ``thou wast alone my solace at home, my charm abroad; \ldots\ the only one on whom rested care of home and of public matters'' (\work{De Excessu} I.20).

Second, he placed himself under intense theological apprenticeship while already bearing the burden of preaching. Augustine later called the elderly Milanese presbyter Simplicianus ``the father of Ambrose in receiving Thy grace'' (\work{Confessions} VIII.2.3), testifying that Simplicianus guided the newly baptized bishop in Scripture and Nicene doctrine. Ambrose read tirelessly in Greek patristic theology---Origen, Athanasius, Didymus, Basil---and admitted candidly to his clergy that he had enjoyed no seminary leisure beforehand:

\begin{quote}
For one is the true Master, Who alone has not learnt, what He taught all; but men learn before they teach, and receive from Him what they may hand on to others. But not even this was the case with me. For I was carried off from the judgment seat, and the garb of office, to enter on the priesthood, and began to teach you, what I myself had not yet learnt. So it happened that I began to teach before I began to learn. Therefore I must learn and teach at the same time, since I had no leisure to learn before. (\work{De Officiis} I.1.3--4)
\end{quote}

Third, he immediately bound the divided Church of Milan back to its exiled Nicene past. Within his first year as bishop, Ambrose sent a delegation of Milanese clerics all the way to Cappadocia to ask Basil of Caesarea for the mortal remains of St.~Dionysius, the Catholic bishop of Milan banished in 355 who had died in eastern exile. Basil's \work{Letter} CXCVII.2 records how reluctantly the Cappadocian faithful parted with the bones they had treasured ``with the honours of a martyr,'' and testifies solemnly to their authenticity: ``One coffin held that honoured corpse. None other lay by his side. \ldots\ Pious were the hands that gave; scrupulously careful were the hands that received. There has been no room for deceit; no room for guile. I bear witness to this.'' By bringing Dionysius home to Milan, Ambrose declared to city and court alike whose episcopal line he continued---not that of the Homoian intruder Auxentius, but the orthodox inheritance of Dionysius, Eustorgius, and Myrocles (\work{Sermon against Auxentius} 18).
